% GENERATED FILE. Edit canonical lessons, then run npm run generate:books.
% canonical-source-hash: fnv1a64:26c27f34d6626e40
% canonical-lessons: BN-C05-bola, BN-C05-ami-bangla-boli, BN-C05-kaj-kora, BN-C05-thaka, BN-C05-practice

\chapter{The First Verbs}
\label{ch:first-verbs}
\hlchaptermodality{5}

\begin{chapteropening}
\textbf{By the end of this chapter:} \emph{I can build short present-tense Bengali sentences with the chapter's first verbs.}

Complete BN-C05-practice, the canonical closing checkpoint, and demonstrate the chapter promise: build short present-tense Bengali sentences with the chapter's first verbs.
\end{chapteropening}

\section[\bn{বলা}]{\bn{বলা} (bôlā) — \textquotedblleft{}to speak,\textquotedblright{} your first full verb}

\label{lesson:BN-C05-bola}



\begin{quote}
So far, \textquotedblleft{}to be.\textquotedblright{} Here is a verb that \emph{does} something — and it reveals how gloriously simple Bengali verbs are.
\end{quote}

\subsection*{The letters in this word}
\textbf{\bn{ব}} (\emph{bô}) + \textbf{\bn{লা}} (\emph{lā}) $\to$ \textbf{\bn{বলা}} (\emph{bôlā}). The dictionary form of the verb.

\begin{cousinweb}
\textbf{\bn{বলা}} (\emph{bôlā}, \textquotedblleft{}to speak, to say\textquotedblright{}) is a native Indo-Aryan verb (the same \emph{bol-} as Hindi \emph{bolnā}, Punjabi \emph{bolṇā}). Its ending \textbf{-\bn{আ}} (\emph{-ā}) is Bengali's infinitive marker. To say \textquotedblleft{}\textbf{I speak},\textquotedblright{} add the present ending \textbf{-\bn{ই}} (\emph{-i}) to the stem: \emph{bôl-} + \emph{-i} $\to$ \textbf{\bn{বলি}} (\emph{bôli}), \textquotedblleft{}I speak.\textquotedblright{}
\end{cousinweb}

\begin{grammarlens}[title={Grammar Lens: the verb changes for person — but \emph{never} for gender}]
Bengali verbs change ending for \textbf{who} speaks: \emph{āmi bôli} (I say), \emph{tumi bôlo} (you say), \emph{se bôle} (he/she says). But — and here is Bengali's great simplification — the verb \textbf{never changes for gender}. There is no \textquotedblleft{}he-form\textquotedblright{} and \textquotedblleft{}she-form,\textquotedblright{} no masculine/feminine at all: \emph{se bôle} is \textquotedblleft{}he says\textquotedblright{} \emph{and} \textquotedblleft{}she says.\textquotedblright{} Where Hindi and Punjabi make even \textquotedblleft{}I speak\textquotedblright{} agree with the speaker's gender (\emph{boltā/boltī}), Bengali dropped grammatical gender entirely. One less thing to track, forever.
\end{grammarlens}

\subsection*{Your turn}
Say these aloud:

\begin{itemize}
\raggedright
  \item \textquotedblleft{}bôlā\textquotedblright{} (to speak)
  \item \textquotedblleft{}I speak\textquotedblright{} — \emph{āmi bôli}
  \item does \textquotedblleft{}he says\textquotedblright{} differ from \textquotedblleft{}she says\textquotedblright{}? (No — \emph{se bôle} for both; Bengali has no gender)
\end{itemize}

\begin{tcolorbox}[breakable,colback=teal!4,colframe=teal!35!black,title={Before you move on}]
What does the Bengali verb change for, and what does it never change for? (For person — \emph{bôli/bôlo/bôle} — but never for gender.)
\end{tcolorbox}
\section[\bn{আমি} \bn{বাংলা} \bn{বলি}]{\bn{আমি} \bn{বাংলা} \bn{বলি} (āmi bānglā bôli) — \textquotedblleft{}I speak Bengali\textquotedblright{}}

\label{lesson:BN-C05-ami-bangla-boli}



\begin{quote}
Your first complete, moving sentence — and the language's own name for itself.
\end{quote}

\subsection*{The letters in this word}
\textbf{\bn{বাংলা}} (\emph{bānglā}): \textbf{\bn{বা}} (\emph{bā}) + \textbf{\bn{ং}} (the \emph{anusvāra} nasal) + \textbf{\bn{লা}} (\emph{lā}).

\begin{cousinweb}
\textbf{\bn{আমি} \bn{বাংলা} \bn{বলি}} = \textbf{\bn{আমি}} (\emph{āmi}, \textquotedblleft{}I\textquotedblright{}) + \textbf{\bn{বাংলা}} (\emph{bānglā}, the object) + \textbf{\bn{বলি}} (\emph{bôli}, \textquotedblleft{}speak\textquotedblright{}) — \textquotedblleft{}I Bengali speak,\textquotedblright{} the verb last as always. The name \textbf{\bn{বাংলা}} (\emph{bānglā}) is the language's own word for itself (English \textquotedblleft{}Bengali\textquotedblright{} comes from it, via \emph{Bangla} $\to$ \emph{Bengal}). It descends from \textbf{\bn{বঙ্গ}} (\emph{bôngo}), the ancient name of the region — so the language is named after its land, the great delta of the Ganges.
\end{cousinweb}

\begin{grammarlens}[title={Grammar Lens: no gender, so the sentence is the same for anyone}]
Because Bengali verbs carry no gender, \textbf{\bn{আমি} \bn{বাংলা} \bn{বলি}} is exactly the same whether a man or a woman says it — unlike Hindi (\emph{maiṁ … boltā/boltī hūṁ}). And because the \emph{-i} ending already means \textquotedblleft{}I,\textquotedblright{} you could even drop \emph{āmi}: \emph{bānglā bôli} is a complete \textquotedblleft{}I speak Bengali.\textquotedblright{}
\end{grammarlens}

\subsection*{Your turn}
Say these aloud:

\begin{itemize}
\raggedright
  \item \textquotedblleft{}āmi bānglā bôli\textquotedblright{}
  \item what \emph{bānglā} is named after (\emph{bôngo}, the ancient region — the Ganges delta)
  \item does the sentence change for a man vs. a woman? (No — Bengali has no gender)
\end{itemize}

\begin{tcolorbox}[breakable,colback=teal!4,colframe=teal!35!black,title={Before you move on}]
Say \textquotedblleft{}I speak Bengali,\textquotedblright{} and give the origin of the name \emph{bānglā}. (\emph{Āmi bānglā bôli}; from \emph{bôngo}, the ancient name of the Bengal region.)
\end{tcolorbox}
\section[\bn{কাজ} \bn{করা}]{\bn{কাজ} \bn{করা} (kāj kôrā) — \textquotedblleft{}to work\textquotedblright{}}

\label{lesson:BN-C05-kaj-kora}



\begin{quote}
The last verb of the chapter — and, like Hindi's \emph{karnā}, it is a \textquotedblleft{}do\textquotedblright{} verb that builds a hundred others.
\end{quote}

\subsection*{The letters in this word}
\textbf{\bn{কাজ}} (\emph{kāj}, \textquotedblleft{}work\textquotedblright{}): \textbf{\bn{কা}} (\emph{kā}) + \textbf{\bn{জ}} (\emph{j}). \textbf{\bn{করা}} (\emph{kôrā}, \textquotedblleft{}do\textquotedblright{}): \textbf{\bn{ক}} (\emph{kô}) + \textbf{\bn{রা}} (\emph{rā}).

\begin{cousinweb}
\textbf{\bn{করা}} (\emph{kôrā}, \textquotedblleft{}to do, to make\textquotedblright{}) is from Sanskrit \textbf{√\bn{কৃ}} (\emph{√kṛ}, \textquotedblleft{}to do\textquotedblright{}), from PIE \textbf{*k\textsuperscript{w}er-} — the same root you met in Chapter 1's \textbf{\bn{নমস্কার}} (\emph{nômoshkar} = \emph{namas} + \emph{kāra}, \textquotedblleft{}the \textbf{making} of a bow\textquotedblright{}) and in \emph{karma}. And \textbf{\bn{কাজ}} (\emph{kāj}, \textquotedblleft{}work\textquotedblright{}) is itself from Sanskrit \emph{kārya} (\textquotedblleft{}work, task,\textquotedblright{} also from √kṛ) — so \textquotedblleft{}work- do\textquotedblright{} is, at root, \textquotedblleft{}\emph{kṛ}-thing \emph{kṛ}-do.\textquotedblright{} Like Hindi, Bengali builds compound verbs as \textbf{noun + \bn{করা}}:

\begin{itemize}
\raggedright
  \item \textbf{\bn{কাজ} \bn{করা}} (\emph{kāj kôrā}) = \textquotedblleft{}work-do\textquotedblright{} = \textquotedblleft{}\textbf{to work}\textquotedblright{}
  \item \textbf{\bn{কথা} \bn{বলা}} (\emph{kôthā bôlā}) = \textquotedblleft{}to converse\textquotedblright{} (using \emph{bôlā} from earlier)
\end{itemize}

So \emph{āmi kāj kôri} = \textquotedblleft{}I work\textquotedblright{} — the \emph{-i} \textquotedblleft{}I\textquotedblright{} ending, no gender.
\end{cousinweb}

\subsection*{Your turn}
Say these aloud:

\begin{itemize}
\raggedright
  \item \textquotedblleft{}kāj kôrā\textquotedblright{} (to work)
  \item \textquotedblleft{}I work\textquotedblright{} — \emph{āmi kāj kôri}
  \item the Chapter-1 word that shares the \emph{√kṛ} root (\emph{nômoshkar})
\end{itemize}

\begin{tcolorbox}[breakable,colback=teal!4,colframe=teal!35!black,title={Before you move on}]
What Sanskrit root do both \emph{kôrā} (\textquotedblleft{}do\textquotedblright{}) and \emph{kāj} (\textquotedblleft{}work\textquotedblright{}) come from, and name a Chapter-1 word built on it. (\emph{√kṛ}, \textquotedblleft{}to do\textquotedblright{}; \emph{nômoshkar} — \textquotedblleft{}the making of a bow.\textquotedblright{})
\end{tcolorbox}
\section[\bn{থাকা}]{\bn{থাকা} (thākā) — \textquotedblleft{}to live, to stay\textquotedblright{}}

\label{lesson:BN-C05-thaka}



\begin{quote}
A second verb, so you can say where you live.
\end{quote}

\subsection*{The letters in this word}
\textbf{\bn{থা}} (\emph{thā}, an aspirated \emph{t} + long \emph{ā}) + \textbf{\bn{কা}} (\emph{kā}) $\to$ \textbf{\bn{থাকা}} (\emph{thākā}).

\begin{cousinweb}
\textbf{\bn{থাকা}} (\emph{thākā}, \textquotedblleft{}to live, to stay, to remain\textquotedblright{}) is from Sanskrit \textbf{\bn{স্থা}} (\emph{sthā}, \textquotedblleft{}to stand, to stay\textquotedblright{}) — the same PIE root \textbf{*steh\textsubscript{2}-} (\textquotedblleft{}to stand\textquotedblright{}) that gives English \textbf{stand}, \textbf{stay}, \textbf{state}, and Latin \emph{stāre}. So Bengali's \textquotedblleft{}to live/stay\textquotedblright{} and English's \textquotedblleft{}to stand/stay\textquotedblright{} are, at the root, one word. \textquotedblleft{}I live in Kolkata\textquotedblright{} is \textbf{\bn{আমি} \bn{কলকাতায়} \bn{থাকি}} (\emph{āmi Kolkātāy thāki}) — stem \emph{thāk-} + the \emph{-i} \textquotedblleft{}I\textquotedblright{} ending, no gender.
\end{cousinweb}

\begin{grammarlens}[title={Grammar Lens: \textquotedblleft{}in\textquotedblright{} comes after the noun}]
\textbf{\bn{কলকাতায়}} (\emph{Kolkātāy}) = \emph{Kolkata} + \textbf{-\bn{য়}} / \textbf{-\bn{তে}} (\emph{-y/-te}, \textquotedblleft{}in, at\textquotedblright{}), glued to the \textbf{end} of the noun. Bengali uses \textbf{post}positions (case-endings) where English uses prepositions: \textquotedblleft{}Kolkata-in I live,\textquotedblright{} the verb last.
\end{grammarlens}

\subsection*{Your turn}
Say these aloud:

\begin{itemize}
\raggedright
  \item \textquotedblleft{}thākā\textquotedblright{}
  \item \textquotedblleft{}I live in Kolkata\textquotedblright{} — \emph{āmi Kolkātāy thāki}
  \item the English cousins of its root \emph{sthā} (\emph{stand, stay, state})
\end{itemize}

\begin{tcolorbox}[breakable,colback=teal!4,colframe=teal!35!black,title={Before you move on}]
What Sanskrit/PIE root is \emph{thākā} from, and its English cousins? (Sanskrit \emph{sthā} / PIE \emph{*steh\textsubscript{2}-} \textquotedblleft{}to stand\textquotedblright{}; English \textbf{stand, stay, state}.)
\end{tcolorbox}
\section[Practice]{Chapter 5 — The first verbs}

\label{lesson:BN-C05-practice}



\begin{quote}
No new words. Three verbs, one gender-free engine — build real sentences about yourself.
\end{quote}

\subsection*{How to answer — build the sentences}
Say these aloud:

\begin{itemize}
\raggedright
  \item \textquotedblleft{}I speak Bengali\textquotedblright{} — \emph{āmi bānglā bôli}
  \item \textquotedblleft{}I live in Kolkata\textquotedblright{} — \emph{āmi Kolkātāy thāki}
  \item \textquotedblleft{}I work\textquotedblright{} — \emph{āmi kāj kôri}
  \item \textquotedblleft{}he speaks\textquotedblright{} and \textquotedblleft{}she speaks\textquotedblright{} — \emph{se bôle} (the same for both!)
\end{itemize}

\begin{grammarlens}[title={Grammar lens — the engine}]
Say these aloud:

\begin{itemize}
\raggedright
  \item the \textquotedblleft{}I\textquotedblright{} ending on every present verb (\emph{-i})
  \item the one thing Bengali verbs never mark (gender)
  \item where the verb sits, and where \textquotedblleft{}in\textquotedblright{} sits (verb last; \emph{-y/-te} after the noun)
\end{itemize}
\end{grammarlens}

\begin{grammarlens}[title={What you've built — roots you now carry}]
\begin{itemize}
\raggedright
  \item \emph{bôlā} \textquotedblleft{}to speak\textquotedblright{} (native; \emph{bol-}, like Hindi/Punjabi)
  \item \emph{thākā} \textquotedblleft{}to live\textquotedblright{} $\leftarrow$ Sanskrit \emph{sthā} $\to$ English \textbf{stand/stay/state}
  \item \emph{kôrā} \textquotedblleft{}to do\textquotedblright{} $\leftarrow$ \emph{√kṛ} — root of \emph{nômoshkar}; \emph{kāj} \textquotedblleft{}work\textquotedblright{} $\leftarrow$ \emph{kārya} (also √kṛ)
  \item \emph{bānglā} $\leftarrow$ \emph{bôngo}, the ancient region — the Ganges delta
\end{itemize}
\end{grammarlens}

\begin{tcolorbox}[breakable,colback=teal!4,colframe=teal!35!black,title={Before you move on}]
In one breath, say your language, your city, and your work in Bengali — and note what stays the same for anyone. (\emph{Āmi bānglā bôli. Āmi … thāki. Āmi kāj kôri} — no gender, so it's the same for a man or a woman.)
\end{tcolorbox}
